% Part: sets-functions-relations
% Chapter: functions
% Section: inverses

\documentclass[../../../include/open-logic-section]{subfiles}

\begin{document}

\olfileid{sfr}{fun}{inv}
\olsection{د تابعو معکوسونه}

\begin{explain}
تابعې د نگاشتونو په توګه ګڼو۔ نو د تابعو په اړه يوه څرګنده
پوښتنه دا ده چې ايا نگاشت «سرچپه» کېدے شي۔ د بېلګې په توګه،
د ځايناستي تابع $f(x) = x + 1$ په دې معنٰی سرچپه کېدے شي چې
تابع $g(y) = y - 1$ د $f$ کړے کار «بېرته اړوي»۔

خو بايد احتياط وکړو۔ که څه هم د~$g$ تعريف يوه تابع $\Int \to \Int$
ټاکي، يوه \emph{تابع} $\Nat \to \Nat$ نۀ ټاکي، ځکه چې
$g(0) \notin \Nat$۔ نو په ساده حالتونو کښې هم دا خبره پوره څرګنده
نۀ ده چې تابع سرچپه کېدے شي که نۀ؛ دا د تعريف په ساحه او هدف
سټ پورې تړلې کېدے شي۔

د تابعې د معکوس مفهوم دا خبره لا کره کوي۔
\end{explain}

\begin{defn}
تابع $g \colon B \to A$ د تابعې $f \colon A \to B$ \emph{معکوس}
ده که د ټولو $x \in A$ او $y \in B$ دپاره $f(g(y)) = y$ او
$g(f(x)) = x$ وي۔
\end{defn}

که $f$ معکوس~$g$ ولري، ډېر وخت د~$g$ پر ځاے $f^{-1}$ ليکو۔

\begin{explain}
اوس به معلومه کړو چې تابعې کله معکوسونه لري۔ د $f\colon A \to B$
د معکوس دپاره يو مناسب نوماند $g\colon B \to A$ دے چې «په دې
تعريف شوے» وي:
\[
g(y) = \text{«هغه يوازينے» $x$ چې $f(x) = y$ وي۔}
\]
خو د «په دې تعريف شوے» (او «هغه يوازينے») شاوخوا احتياطي نښې
اشاره کوي چې دا تعريف نۀ دے۔ لږ تر لږه، تل او په بشپړ عموميت
سره کار نۀ ورکوي۔ ځکه د دې دپاره چې دا تعريف يوه تابع وټاکي،
بايد يو او يوازې يو~$x$ وي چې $f(x) = y$ وي---د~$g$ وت بايد په
يوازيني ډول ټاکل شوے وي۔ سربېره پر دې، د هر $y \in B$ دپاره
بايد ټاکل شوے وي۔ که داسې $x_1$ او $x_2 \in A$ وي چې
$x_1 \neq x_2$ خو $f(x_1) = f(x_2)$ وي، نو د
$y = f(x_1) = f(x_2)$ دپاره به $g(y)$ په يوازيني ډول نۀ وي
ټاکل شوے۔ او که هېڅ داسې~$x$ نۀ وي چې $f(x) = y$ وي، نو $g(y)$
به بېخي نۀ وي ټاکل شوے۔ په بله وينا، د~$g$ د تعريف کېدو دپاره
بايد~$f$ هم !!{injective} وي او هم !!{surjective}۔

راځئ ګام په ګام لاړ شو۔ پوښتنه به په دوو برخو ووېشو: د ورکړل
شوې تابعې~$f\colon A \to B$ دپاره، کله داسې تابع $g\colon B \to A$
شته چې $g(f(x)) = x$ وي؟ داسې~$g$ د $f$ کړے کار «بېرته اړوي»
او د~$f$ \emph{کيڼ معکوس} بلل کېږي۔ دويمه دا چې کله داسې تابع
$h\colon B \to A$ شته چې $f(h(y)) = y$ وي؟ داسې~$h$ د~$f$
\emph{ښي معکوس} بلل کېږي---$f$ د~$h$ کړے کار «بېرته اړوي»۔
\end{explain}

\begin{prop}
که د تعريف ساحه ناتشه وي او $f\colon A \to B$ يوه !!{injective} تابع وي، نو د~$f$ يو
\emph{کيڼ معکوس}~$g\colon B \to A$ شته چې د ټولو $x \in A$
دپاره $g(f(x)) = x$ وي۔
\textbf{د ژباړې د نسخې سمون (OLFUN-001):} په انګرېزي سرچينه کښې
د ناتشې تعريف ساحې شرط نۀ و؛ خو لاندې ثبوت يو غړے ترې ټاکي۔
\end{prop}

\begin{proof}
فرض کړئ $f\colon A \to B$ يوه !!{injective} تابع ده۔ يو $y \in B$
په پام کښې ونيسئ۔ که $y \in \ran{f}$ وي، نو داسې $x \in A$ شته
چې $f(x) = y$ وي۔ ځکه چې $f$ !!{injective} ده، يوازې يو داسې
$x \in A$ شته۔ بيا داسې تعريفولے شو: $g(y) = x$؛ يعنې $g(y)$
«هغه يوازينے» $x \in A$ دے چې $f(x) = y$ وي۔ که
$y \notin \ran{f}$ وي، له هر يوه~$a \in A$ سره ئې تړلے شو۔ نو
يو $a \in A$ ټاکلے او $g\colon B \to A$ داسې تعريفولے شو:
\[
g(y) = \begin{cases}
    x & \text{که $f(x) = y$ وي}\\
    a & \text{که $y \notin \ran{f}$ وي۔}
\end{cases}
\]
دا د ټولو $y \in B$ دپاره تعريف شوې ده، ځکه چې د هر داسې
$y \in \ran{f}$ دپاره په دقيقه توګه يو $x \in A$ شته چې
$f(x) = y$ وي۔ د تعريف له مخې، که $y = f(x)$ وي، نو $g(y) = x$
وي؛ يعنې $g(f(x)) = x$۔
\end{proof}

\begin{prob}
وښيئ چې که $f\colon A \to B$ کيڼ معکوس~$g$ ولري، نو~$f$
!!{injective} ده۔
\end{prob}

\begin{prop}
    که $f\colon A \to B$ يوه !!{surjective} تابع وي، نو د~$f$
    يو \emph{ښي معکوس}~$h\colon B \to A$ شته چې د ټولو~$y \in B$
    دپاره $f(h(y)) = y$ وي۔
\end{prop}

\begin{proof}
فرض کړئ $f\colon A \to B$ يوه !!{surjective} تابع ده۔ يو $y \in B$
په پام کښې ونيسئ۔ ځکه چې~$f$ !!{surjective} ده، داسې $x_y \in A$
شته چې $f(x_y) = y$ وي۔ بيا داسې تعريفولے شو: $h(y) = x_y$؛
يعنې د هر $y \in B$ دپاره يو داسې $x \in A$ ټاکو چې $f(x) = y$
وي؛ ځکه چې~$f$ !!{surjective} ده، تل لږ تر لږه يو د ټاکلو دپاره
شته۔\footnote{ځکه چې $f$ !!{surjective} ده، د هر~$y \in B$
دپاره سټ $\Setabs{x}{f(x) = y}$ تش نۀ دے۔ زموږ د~$h$ تعريف غواړي
چې له دې هر سټ څخه يو يوازينے $x$ وټاکو۔ دا چې دا کار تل ممکن
دے، په حقيقت کښې څرګنده خبره نۀ ده---د دې ټاکنو امکان په ساده
ډول د يوه بديهي اصل په توګه فرض کېږي۔ په بله وينا، دا قضيه
د ټاکنې په نامه بديهي اصل فرض کوي؛ دا مسئله به
\oliflabeldef{sth:choice::chap}{په \olref[sth][choice][]{chap} کښې
بيا وڅېړو}{په لنډه پرېږدو}۔ خو په ډېرو ځانګړو حالتونو کښې، لکه
چې $A = \Nat$ وي يا متناهي وي، يا $f$ يوه !!{bijective} تابع وي،
د ټاکنې بديهي اصل ته اړتيا نۀ وي۔ (په هغه ځانګړي حالت کښې چې
$f$ !!{bijective} وي، د هر $y \in B$ دپاره سټ
$\Setabs{x}{f(x) = y}$ په دقيقه توګه يو !!{element} لري؛ نو د
انتخاب دپاره څو بديلونه نشته۔)}
د تعريف له مخې، که $x = h(y)$ وي، نو $f(x) = y$ وي؛ يعنې
د هر $y \in B$ دپاره $f(h(y)) = y$۔
\end{proof}

\begin{prob}
وښيئ چې که $f\colon A \to B$ ښي معکوس~$h$ ولري، نو~$f$
!!{surjective} ده۔
\end{prob}

\begin{explain}
  د مخکيني ثبوت د مفکورو په يوځاے کولو اوس دا ترلاسه کوو چې
  هره !!{bijection} معکوس لري؛ يعنې يوه واحده تابع شته چې
  هم د~$f$ کيڼ معکوس دے او هم ښي معکوس۔
\end{explain}

\begin{prop}\ollabel{prop:bijection-inverse}
که $f\colon A \to B$ يوه !!{bijective} تابع وي، داسې تابع
$f^{-1}\colon B \to A$ شته چې د ټولو $x \in A$ دپاره
$f^{-1}(f(x)) = x$ او د ټولو $y \in B$ دپاره $f(f^{-1}(y)) = y$ وي۔
\end{prop}

\begin{proof}
تمرين۔
\end{proof}

\begin{prob}
\olref[sfr][fun][inv]{prop:bijection-inverse} ثابته کړئ۔ بايد
$f^{-1}$ تعريف کړئ، وښيئ چې تابع ده، او وښيئ چې د~$f$ معکوس
ده؛ يعنې د ټولو $x \in A$ او $y \in B$ دپاره $f^{-1}(f(x)) = x$
او $f(f^{-1}(y)) = y$ وي۔
\end{prob}

\begin{explain}
د معکوسونو د ترلاسه کولو يوه لږه زياته عمومي لار هم شته۔ په
\olref[kin]{sec} کښې مو وليدل چې هره تابع $f$ د ټولو $x \in A$
دپاره د $f'(x) = f(x)$ په ټاکلو سره يوه !!a{surjection}
$f' \colon A \to \ran{f}$ رامنځته کوي۔ په ښکاره، که~$f$
!!{injective} وي، نو~$f'$ !!{bijective} ده، او له دې کبله د
\olref{prop:bijection-inverse} له مخې يوازينے معکوس لري۔ د نښو
په کارولو کښې په يو ډېر واړه تسامح سره، کله کله د $f'$ معکوس
يوازې «د~$f$ معکوس» بولو۔
\end{explain}

\begin{prop}\ollabel{prop:left-right}%
  وښيئ چې که $f\colon A \to B$ کيڼ معکوس~$g$ او ښي معکوس~$h$
  ولري، نو $h = g$۔
\end{prop}

\begin{proof}
  تمرين۔
\end{proof}

\begin{prob}
  \olref[sfr][fun][inv]{prop:left-right} ثابته کړئ۔
\end{prob}

\begin{prop}\ollabel{prop:inverse-unique}
هره تابع~$f$ تر يوه زيات معکوس نۀ لري۔
\end{prop}

\begin{proof}
  فرض کړئ $g$ او $h$ دواړه د~$f$ معکوسونه دي۔ نو په ځانګړي ډول
  $g$ د~$f$ کيڼ معکوس او $h$ ښي معکوس دے۔ د
  \olref{prop:left-right} له مخې، $g = h$۔
\end{proof}

\end{document}
